\section{本章小结}
\label{sec:dcc-summary}

本章研究有限字母表上、编码器不使用反馈的离散无记忆信道。信道由转移矩阵$\bm{W}$给出，发送端能选择的是输入分布与码本；可靠传输的含义是码长增长时错误概率趋于零，速率按比特每次信道使用计量。

\paragraph{主要结论}
\begin{itemize}
 \item 容量公式：$C=\max_qI(q)$。$I(q)$是单纯形上的连续凹函数，最大值存在，最优输出分布唯一；最优输入由条件~\eqref{eq:dcc-kkt}刻画：被使用的输入的散度$D_{\mathrm{KL}}(W_x\Vert p^{*})$都等于$C$，未使用的不超过$C$。对任意输出分布$r$，$\max_xD_{\mathrm{KL}}(W_x\Vert r)$是容量的上界。

 \item 块的信息：无记忆与无反馈保证$I(X^n;Y^n)\leq nC$，独立使用最优输入达到等号。

 \item 编码定理：码率$R<C$时，随机码本加联合典型解码使平均错误趋于零，删去一半消息后最大错误也趋于零；码率$R>C$时，Fano不等式与数据处理给出$P_{\mathrm{e}}\geq1-C/R-1/(nR)$，错误不能趋于零。定理对$R=C$及有限块长不作承诺。

 \item 二元信道：二元对称信道$C=1-h_2(p)$，二元擦除信道$C=1-\varepsilon$；均匀输入最优来自交换对称性，Z信道的最优输入是$0.4$。已知错误位置的擦除比未知位置的翻转损失少，差值$h_2(p)-p$是定位错误的代价。

 \item 成本约束：逐条码字的预算下，容量是容量—成本函数$C(\Gamma)$，它单调、凹，切线斜率是预算的边际价值；可达性需要为预算留出余量，预算等于最小成本时改用受限字母表。

 \item 数值计算：Blahut–Arimoto迭代按$2^{D_{\mathrm{KL}}(W_x\Vert p_t)}$乘性地重新加权输入，每轮给出严格的上下界，下界误差不超过$D_{\mathrm{KL}}(q^{*}\Vert q_0)/(T+1)$；初始分布须有完整支撑，停止依据是上下界之差而不是迭代的变化量。
\end{itemize}

\paragraph{对比}
本章有几组容易混淆的概念，表~\ref{tab:dcc-summary-contrast}把它们并列。

\begin{BookInlineTable}
 \centering
 \caption{本章几组相邻概念的区别。}
 \label{tab:dcc-summary-contrast}
 \begin{tabular}{p{18mm}p{39mm}p{39mm}}
  \toprule
  对象 & 前者 & 后者 \\
  \midrule
  两个容量 & $C_{\mathrm{op}}$：可靠码率的上确界 & $C$：单次互信息的最大值 \\
  两种错误 & 平均错误：均匀消息下的错误率 & 最大错误：每条消息的保证 \\
  定理两向 & 可达性：存在好码，不给构造 & 逆定理：对一切码成立的限制 \\
  两种损坏 & 翻转：位置未知，$1-h_2(p)$ & 擦除：位置可见，$1-\varepsilon$ \\
  两种成本 & 期望成本：输入分布的约束 & 逐条预算：每条码字的约束 \\
  两种停止 & 迭代变化小：不是证书 & $U_t-L_t$小：严格的误差保证 \\
  \bottomrule
 \end{tabular}
\end{BookInlineTable}

容量回答“最多能多快”，没有回答“给定块长需要多少冗余”和“怎样以可承受的计算量编码与解码”。随机码本不可存储，典型性解码要搜索指数多个码字；有限长度的纠错构造与性能比较由第~\ref{chap:channel-codes}章展开。下一章先把输入字母从有限集合换成受功率限制的实数，研究高斯噪声下的容量，见第~\ref{chap:gaussian-channel-capacity}章。
